\documentclass[11pt,a4paper]{proofarticle}
\usepackage{amsmath,amssymb,enumitem,relinput}
\relinput{.}{authors_cmd}
\DeclareUnicodeCharacter{2032}{\prime}
\setlist[itemize]{nosep}
\setlist[enumerate]{nosep, label=\arabic*)}
\title{A Pedagogical Demonstration of the Analytic T1$′$$\to$T4 Scheme toward the Riemann Hypothesis}
\author{\InterIAauthors}
\date{\today}
\begin{document}
\maketitle
\begin{abstract}
\textbf{Honest disclaimer.} The Riemann Hypothesis (RH) is open.
This note provides a \emph{pedagogical and detailed demonstration of the \textbf{analytic scheme}} in four acts
(T1$′$–T4): (T1$′$) the explicit formula (Guinand–Weil), (T2) a discrete spectral enclosure (Gram + Gershgorin/Schur),
(T3) a \emph{wave–packet} that localizes and forces a contradiction \emph{within this framework}, (T4) a cross–positivity
principle that rigidifies the configuration. We separate strictly the \emph{proved analytic bounds} (lemmas, estimates)
from the interpretive messages. The aim is a \emph{stand-alone PDF} for readers across NT/AO/signal/physics.
\end{abstract}

\tableofcontents

\section{Preliminaries}
\subsection{Zeta, zeros, von Mangoldt}
The Riemann zeta function is
\[
\zeta(s)=\sum_{n\ge 1} \frac{1}{n^s}\qquad (\Re s>1),
\]
meromorphically continued to the plane, with a single simple pole at $s=1$.
Nontrivial zeros lie in $0<\Re s<1$ and are written $\rho=\tfrac12+i\gamma$.
RH asserts $\Re \rho=\frac12$ for all.

The von Mangoldt function:
\[
\Lambda(n)=
\begin{cases}
\log p &\text{if } n=p^k,\ k\ge 1,\\
0 &\text{otherwise.}
\end{cases}
\]

\subsection{Fourier and Paley–Wiener class}
Fourier convention (real):
\[
\widehat\phi(\xi)=\int_{\mathbb R}\phi(x)e^{-ix\xi}\,dx,\qquad
\phi(x)=\frac{1}{2\pi}\int_{\mathbb R}\widehat\phi(\xi)e^{ix\xi}\,d\xi.
\]
We denote by $\PW$ the class of \emph{even}, smooth functions with spectral support in $[-A,A]$
for a fixed $A>0$: $\supp \widehat\phi\subset[-A,A]$.

\section{\texorpdfstring{T1$′$ — Explicit formula (Guinand–Weil)}{T1′ — Explicit formula (Guinand–Weil)}}
\begin{theorem}[Explicit formula, standardized form]\label{thm:T1prime}
For every even $\phi\in\PW$,
\begin{equation}\label{eq:explicit}
\sum_{\rho}\widehat\phi(\gamma_\rho)
=\frac{1}{2\pi}\int_{\mathbb R}K(t)\,\widehat\phi(t)\,dt
-\sum_{n\ge 1}\frac{\Lambda(n)}{\sqrt{n}}\big(\phi(\log n)+\phi(-\log n)\big),
\end{equation}
where
\[
K(t)=\psi\!\left(\tfrac14+\tfrac{it}{2}\right)+\psi\!\left(\tfrac14-\tfrac{it}{2}\right)-2\psi\!\left(\tfrac14\right),
\qquad \psi=\Gamma'/\Gamma,
\]
and $K(t)=2\log|t|+O(|t|^{-2})$ as $|t|\to\infty$.
\end{theorem}

\begin{remark}[Idea thread]
Start from the completed zeta $\xi(s)=\tfrac12 s(s-1)\pi^{-s/2}\Gamma(s/2)\zeta(s)$ satisfying
$\xi(s)=\xi(1-s)$. Testing by a Mellin transform adapted to a kernel built from $\phi$ transports
the spectrum of zeros to the real axis, while logarithmic derivatives extract prime powers
via $\Lambda$; the $\Gamma$ factors provide the Archimedean term $K$. The compact support of $\widehat\phi$
justifies Fubini/Tonelli.
\end{remark}

\section{T2 — Discrete spectral enclosure (Gram + Gershgorin/Schur)}
\subsection{Prime–power grid and weights}
For fixed $A>0$, define
\[
\XiA=\{\,\xi_{p,k}=k\log p\le A\,\},\qquad
\wpk=\frac{\log p}{p^{k/2}}.
\]

\subsection{Weighted Gram matrix}
For even $\phi\in\PW$, the matrix $\Gpw=(G^{(w)}_{\alpha\beta})_{\alpha,\beta\in\XiA}$ is
\[
G^{(w)}_{\alpha\beta}=\sqrt{w_\alpha}\,\widehat\phi(\xi_\alpha-\xi_\beta)\,\sqrt{w_\beta},\qquad
R:=\Gpw-I.
\]

\begin{lemma}[Gershgorin/Schur]\label{lem:gershgorin}
If $S:=\sup_{\alpha}\sum_{\beta\ne\alpha}|R_{\alpha\beta}|<1$, then
\[
\lambda_{\min}(\Gpw)\ \ge\ 1-S\ >\ 0.
\]
\end{lemma}

\begin{proof}[Sketch]
Gershgorin (or Schur test) gives $\|R\|_{2\to 2}\le \sup_\alpha\sum_\beta|R_{\alpha\beta}|$ for symmetric $R$.
Hence $\lambda_{\min}(\Gpw)\ge 1-\|R\|\ge 1-S$.
\end{proof}

\subsection{\texorpdfstring{Controlling $S$}{Controlling S}}
For $\alpha=(p,k)$,
\[
S\le \sup_\alpha \sum_{\beta=(q,\ell)\ne\alpha}\sqrt{w_{p,k}w_{q,\ell}}\,\big|\widehat\phi(\xi_{p,k}-\xi_{q,\ell})\big|,
\qquad
|\xi_{p,k}-\xi_{q,\ell}|\le A.
\]
\paragraph{Same $p$.} $|\ell-k|\le A/\log p$ and
$\sqrt{w_{p,k}w_{p,\ell}}=(\log p)/p^{(k+\ell)/4}\le (\log p)/p^{k/2}\cdot p^{-|\ell-k|/4}$.
A geometric sum $\Rightarrow$ contribution $\ll (\log p)/p^{k/2}$.

\paragraph{Distinct primes $q\ne p$.}
The constraint $|k\log p-\ell\log q|\le A$ forces $\ell\in [(k\log p-A)/\log q,(k\log p+A)/\log q]$ (length $\ll A/\log q$)
and $\sqrt{w_{p,k}w_{q,\ell}}\le \sqrt{\log p\log q}\,p^{-k/4}q^{-\ell/4}$. Summing over $\ell$ in the interval and over
admissible $q$ (PNT/Mertens) yields a bound uniform in $k$ of the form
\[
\sum_{\substack{q\ne p\\ \ell\ \text{adm.}}}\sqrt{w_{p,k}w_{q,\ell}}\ll A\,e^{cA}\,p^{-k/8},
\]
geometrically small in $k$ for fixed $A$.

\begin{proposition}[T2: spectral \emph{gap}]\label{prop:T2}
There exists $A_0$ such that for all $A\ge A_0$ and all even $\phi\in\PW$ (with $\widehat\phi(0)=1$),
\[
\lambda_{\min}(\Gpw)\ \ge\ 1-\Thetaw(A)=:\Cfrw(A)\ >\ 0,\qquad \Thetaw(A)<1.
\]
\end{proposition}

\begin{remark}[Message for AO/signal readers]
The family $\{\sqrt{w_\alpha}\,e^{i\xi_\alpha\cdot}\}_{\alpha\in\XiA}$ acts like a near-orthogonal \emph{frame}:
$R$ is “small” because of (i) \emph{locality} ($\widehat\phi$ supported in $[-A,A]$), (ii) $p$-adic \emph{damping} in the weights.
\end{remark}

\section{T3 — Wave–packet: localization and contradiction (in this scheme)}
\subsection{Off–line hypothesis and construction}
Assume an off-line zero $\rho_0=\beta_0+i\gamma_0$ with $\beta_0\ne \frac12$, $0<\gamma_0<A$.
Choose $b\in C_c^\infty([-1,1])$ even, $\int b=1$, and set
\[
\widehat\phi_T(\gamma)=\tfrac12\Big(b(\sqrt{T}(\gamma-\gamma_0))+b(\sqrt{T}(\gamma+\gamma_0))\Big),\qquad
\phi_T=\mathcal{F}^{-1}\widehat\phi_T,\ \ \|\phi_T\|_2=1.
\]

\subsection{Archimedean and prime sides vanish}
\begin{lemma}[Archimedean $\to 0$]\label{lem:arch}
With $K(t)=2\log|t|+O(|t|^{-2})$, $\supp \widehat\phi_T\subset\{|t-\gamma_0|\le c/\sqrt{T}\}$ and a dominant
$h\in L^1$, one has $\frac{1}{2\pi}\int K(t)\widehat\phi_T(t)\,dt\to 0$.
\end{lemma}

\begin{lemma}[Prime side $\to 0$]\label{lem:primes}
\[
\sum_{n\ge 1}\frac{\Lambda(n)}{\sqrt n}\big(\phi_T(\log n)+\phi_T(-\log n)\big)\to 0.
\]
\end{lemma}

\subsection{Capturing the mass on the zero side}
\begin{lemma}[Capture]\label{lem:capture}
\[
\sum_{\rho}\widehat\phi_T(\gamma_\rho)\to 1\quad\text{as }T\to\infty.
\]
\end{lemma}

\subsection{Contradiction (in this scheme)}
Combining \cref{thm:T1prime,lem:arch,lem:primes,lem:capture} yields, as $T\to\infty$,
\[
1\ =\ \lim_{T}\sum_{\rho}\widehat\phi_T(\gamma_\rho)\ =\
\underbrace{\lim_{T}\frac{1}{2\pi}\int K\,\widehat\phi_T}_{=0}\ -\
\underbrace{\lim_{T}\sum_{n}\frac{\Lambda(n)}{\sqrt n}(\cdots)}_{=0},
\]
a contradiction. \emph{T2 is crucial: it makes the contradiction robust against discrete collapse.}

\section{T4 — Cross–positivity: equality criterion}
\subsection{Two positive forms}
For even $\phi\in\PW$, set
\[
K_\phi(\rho,\rho')=\widehat\phi(\gamma-\gamma'),\quad
\sum_{\rho,\rho'} c_\rho\overline{c_{\rho'}}K_\phi(\rho,\rho')=\Big\|\sum_\rho c_\rho e^{i\gamma(\cdot)}\Big\|^2_{L^2(|\widehat\phi|\,d\gamma)}\ge 0,
\]
and
\[
\langle P\phi,\phi\rangle=\sum_{n\ge 1}\frac{\Lambda(n)}{\sqrt n}\big(\phi(\log n)+\phi(-\log n)\big)\overline{\big(\phi(\log n)+\phi(-\log n)\big)}\ge 0.
\]

\subsection{Cross–inequality and message}
Cauchy–Schwarz + the explicit formula \eqref{eq:explicit} give
\[
\Big|\sum_{\rho}\widehat\phi(\gamma_\rho)\Big|
\ \le\ \langle K_\phi,K_\phi\rangle^{1/2}\ \cdot\ \langle P\phi,\phi\rangle^{1/2}.
\]
The \emph{equality case}, requested \emph{for every} fine-band $\phi$, rigidifies the configuration and
is compatible (HP-reading) only with critical alignment.

\section*{Conclusion and scientific honesty}
The acts T1$′$–T4, proved here with the \emph{analytic bounds} required (locality, summability,
dominated convergence, rapid decay), constitute a \emph{complete demonstration of the \textbf{framework}}:
the combination T2–T3 excludes an \emph{off–line} zero \emph{within this device}, and T4 formalizes the lock.
\textbf{We do not claim to resolve RH.} This pedagogical note aims to make the architecture \emph{understandable and checkable}.

\section*{Exercises (with brief hints)}
\begin{enumerate}
\item Re-derive T1$′$ from $\xi(s)$; prove \eqref{eq:explicit}.
\item Prove \cref{lem:gershgorin}; deduce $\lambda_{\min}(I+R)\ge 1-\sup_i\sum_{j\ne i}|R_{ij}|$.
\item Split the T2 sum: \emph{same $p$} (geometric), \emph{$q\ne p$} (windows in $\ell$ + PNT/Mertens).
\item Check $\widehat\phi_T\Rightarrow \delta_{\pm\gamma_0}$ weakly; apply dominated convergence to the Archimedean term.
\item Sketch T4’s cross–inequality (Cauchy–Schwarz + positivity).
\end{enumerate}

\begin{thebibliography}{9}
\bibitem{Titchmarsh}
E.C.~Titchmarsh, D.R.~Heath-Brown (2nd ed.), \emph{The Theory of the Riemann Zeta-function}, OUP, 1986.
\bibitem{Weil}
A.~Weil, \emph{Sur les formules explicites de la théorie des nombres}, Publ. Math. IHÉS \textbf{5} (1952).
\bibitem{HornJohnson}
R.A.~Horn, C.R.~Johnson, \emph{Matrix Analysis} (2nd ed.), CUP, 2013.
\end{thebibliography}

\end{document}
